\chapter{基于人工的知识获取}
\label{chap:knowledge-acquisition-human}

假设我们已经收集了一批产品使用文档，需要从中找出故障报告，并取得支持判断的原句。人工预算只够阅读其中一部分，因此不能把每篇文档都交给专家逐项检查。即使决定了“请人来标”，仍有几个问题没有解决：问整篇文档的类别还是问某句的含义，要求直接作答还是核对建议，先处理哪些文档，由谁判断，以及争议样本应追加几次检查。只增加答案条数可能留下没有依据的标签，只询问模型最困惑的材料则可能反复遇到缺页或无法判定的文档。

本章围绕这项文档任务讨论有限人工资源怎样使用。材料中既有明确的常见故障，也有稀有故障、专业术语和仅凭可见上下文无法判断的情形。这里要辨认的是文档是否报告了故障，不是仅凭作者的话就证明设备确实存在故障；若任务需要后者，还须提供检测等相应证据。类别判断与证据片段定位是主要输出，图像关键点和对象边界只用于解释局部标注增加了什么信息，不把描绘几何边界的成本套到阅读文本上。

我们可以按决策对象把这个方案分开考察。任务设计先确定答案的含义，交互方式再确定人怎样提供答案；在可执行的任务上，样本选择决定询问哪些文档，标注者选择决定把材料交给谁，预算分配决定各项行动执行多少次。这五项相互制约，但不能用一个总分代替。人工记录沿用第~\ref{chap:knowledge-acquisition-overview}章的接口：执行状态说明操作进展，来源动作区分作答与弃权，语义状态说明内容是否确定，验证状态说明是否经核验。本章补充任务依据、操作和成本，后续章节再处理答案融合与独立质量审计。

\section{标注任务的设计}
\label{sec:knowledge-acquisition-human-task-design}

把一篇文档交给标注者，并说“帮助提高故障识别能力”，还不能确定应当得到什么记录。人需要知道检查哪个对象、可以依据哪些材料，以及什么情形允许给出确定答案。任务设计就是把下游需求转成这些可执行的判断，同时规定材料或规则不足时怎样停止作答。任务定义、规范和界面还需要经过小规模试行；试行若暴露缺口，应改动相应环节，再检查原有答案是否仍成立。

\subsection{标注目标的确定}
\label{sec:knowledge-acquisition-human-task-objective}

先把贯穿任务写成一条可以交付的请求：检查一篇文档，只使用正文及给定附件，回答“作者是否报告了产品故障”，并保存支持答案的原句或无法判断的原因。待判断对象、允许依据、提问、回答空间和补充记录构成了任务的五个要素。界面可以提供“是、否、未知、不适用”，但训练用的类别仍是“是、否”；“未知”说明内容尚未确定，“不适用”说明对象不在任务范围。人完成检查后因依据不足而不输出判断时，来源动作记为“弃权”。只保存一个“是”而丢掉任务版本和证据范围，就无法分辨它回答的是作者亲历、所有转述，还是经过检测确认的故障。

在这个请求中，文档类别和证据片段承担不同作用。类别概括整篇材料，因此要阅读足以判断事件与说话者的上下文；片段则指出哪些文字支持这个类别，还要确定起止位置。几句文字共同支持结论时，类别可能已经确定，却没有唯一证据句；只提到异常现象的一句，也可能不足以判断整篇文档是不是故障报告。若训练确实需要两种输出，可以先判类别，再对正类定位证据，用稳定的对象标识关联两条记录。没有选中片段只说明证据定位尚未完成或没有找到合适片段，不能直接推出负类。

如果再要求标注者写理由，就要确定这段文字供谁使用。争议复核需要引用位置、采用的条款和缺失信息，以便另一人重做判断；训练文本生成模型还需要限定事实范围、长度和可接受表述，以免把任意发挥当成目标答案。作答后写出的理由可以支持复核，却未必重现人实际形成判断的心理过程。没有额外验证时，不能把这类事后说明当作忠实的因果解释。

一篇文档可能足以判断故障及其部件，却没有严重程度的依据。把“是否故障、故障部件、严重程度、证据片段”拆成四项，就能保留已经确定的两项，并把需要专业知识的一项单独转交。拆分也有代价：材料可能被反复打开，各项结果还要重新关联到同一文档和对象。因此我们应看成本发生在哪里。阅读全文最费时，拆分可能增加重复阅读；精细操作或专业裁决最费时，拆分才可能避免让所有样本都经历最昂贵的步骤。

同样是没有确定答案，下一步行动可能完全不同。附件缺失属于“信息不足”，需要补材料；看到了专业报告却无法理解，属于“缺少知识”，可以转给具备相应能力的人；规则没有说明某种边界情形怎样判断，属于“规范歧义”，应先交给规范负责人处理。三种原因下，内容的语义状态都先保留为“未知”，完成检查而不作判断的来源动作都是“弃权”。分开记录原因，才能避免对缺材料的文档不断换人，或对规则不清的任务不断追加投票。

% 小节入口连同首段长引文预留，避免边注伸入页底留白。
\Needspace{95mm}
\subsection{标注规范的制定}
\label{sec:knowledge-acquisition-human-guideline}

两位标注者都理解“故障”这个词，却可能分别纳入用户误用、转述投诉或尚未确认的怀疑。要使他们回答同一个问题，需要\term{标注规范}（annotation guideline），也就是能在不同文档上反复应用的判断规则。最小规范应说明目标对象、类别定义、纳入与排除条件、对象边界、判断顺序，以及冲突或证据不足时交给谁处理。类别名称只能指向规则，不能代替规则本身。大规模检索标注还会受到疲劳、展示位置和先见建议所形成的锚定影响；为解释结果差异，规范应连同界面顺序、批次长度和休息条件一起记录\scite{knowledge-acquisition-human-thomas2022crowd}

我们可以用“论坛里有人说电池坏了”检验这条规则。它含有“电池坏了”，但故障主张来自被转述的人。若任务限定作者亲历，就要检查说话者与事件归属，并将该句排除；若任务收集所有被提及的故障主张，该句又应纳入。同一关键词得到不同答案，是因为对象与纳入条件改变了，不是因为标注者识别关键词的能力不同。正例说明如何命中条件，反例说明排除依据，边界例则用于补清证据缺失或归属冲突的处理规则。

判断文档是否包含某项可核查事实，与判断某种回答是否可接受，所需规范也不同。前者应给出核查依据和允许查阅的外部资料；后者要固定评价维度、观察条件和量表含义，使不同人的看法可以按同一口径记录，而不是强行指定唯一事实答案。随机化展示顺序或约束流程，有助于控制相应的已知影响，但不能自动消除主观差异或未观察到的偏差。

表~\ref{tab:knowledge-acquisition-human-guideline-cases}把任务固定为作者亲历的故障报告，让我们逐行追踪“对象归属和证据是否满足规则”。第二行的“否”是否定作者亲历，不是否定电池可能损坏；后两行保留未知，分别送去补材料或澄清规范。表中输入都是教学构造，只说明规则怎样执行，不表示四行已经覆盖全部语言边界。

\begin{table*}[htbp]
  \centering
  \small
  \begingroup
  \setlength{\tabcolsep}{4pt}
  \renewcommand{\arraystretch}{1.15}
  \begin{tabularx}{\linewidth}{@{}>{\RaggedRight\setlength{\parindent}{0pt}\arraybackslash}p{29mm}>{\RaggedRight\setlength{\parindent}{0pt}\arraybackslash}p{27mm}>{\RaggedRight\setlength{\parindent}{0pt}\arraybackslash}p{34mm}>{\RaggedRight\setlength{\parindent}{0pt}\arraybackslash}p{14mm}>{\RaggedRight\setlength{\parindent}{0pt}\arraybackslash}X@{}}
    \toprule
    输入示例 & 判断对象 & 适用条款 & 输出 & 证据与去向 \\
    \midrule
    “我升级后无法开机。” & 作者本人的设备经历 & 故障明确且归属于作者 & 是 & 保留整句及文档标识 \\
    “论坛里有人说电池坏了。” & 作者本人的设备经历 & 转述不计入作者亲历 & 否 & 保留说话者线索 \\
    “升级后情况更糟，详见附件。”但附件缺失 & 当前文档的故障判断 & 所需证据不在可见材料中 & 未知 & 信息不足：记录缺少附件，转入补材料队列 \\
    设备与外接配件并述，代词指向不明且规范未给判定顺序 & 哪个对象发生故障 & 对象归属须唯一或可核查；判定顺序缺失 & 未知 & 规范歧义：反馈规范负责人 \\
    \bottomrule
  \end{tabularx}
  \endgroup
  \caption{标注规范怎样落实为可执行判断。四行都判断作者亲历的故障报告：“是、否”针对这一任务，“未知”保留尚未确定的内容。规则同时决定保留哪项证据和交给谁继续处理；教学样例不表示已覆盖全部边界。}
  \label{tab:knowledge-acquisition-human-guideline-cases}
\end{table*}

规则写清后，还须让界面和培训使用同一套规则。三者应共享版本标识，界面选项对应规范的回答空间，培训反馈指出具体违反了哪条规则。规范从版本$v$更新到$v+1$时，要列出改变答案含义的条款及受影响样本，再决定重标、复核还是保留旧版本。若规则语义已经改变，却只替换说明文字并直接合并新旧标签，同一个类别名称就可能包含两套标准。

\subsection{标注任务的试标与调整}
\label{sec:knowledge-acquisition-human-pilot}

即使规范在纸面上完整，正式人员也可能找不到附件、误解选项或耗费过长时间。扩大规模前，应让拟参与工作的人员按完整流程处理一小批材料，这就是\term{人工试标}（human pilot annotation），即由人执行的小规模试运行。样本要同时包含典型情形、预期难例和已知边界例，并保存答案、证据、弃权原因、阅读与操作时间以及疑问。它回答任务能否执行、哪里受阻和怎样估计成本；数据集整体质量及统计不确定性由第~\ref{chap:knowledge-quality-assessment-correction}章处理。“试标”还可以由模型执行，第~\ref{sec:knowledge-acquisition-model-pilot-validation}节讨论那种不同的执行主体。

试标出现不同答案时，我们应先检查人是否读到同样的材料、采用同一条规则。多人对条款作不同解释，可能需要改规范；证据根本没有显示，增加培训也不能补上；规则与信息都充分，错误执行才直接指向培训、界面或人员准入。难例可以讨论，但要先保存各人的独立初始答案，否则讨论后的共识会掩盖分歧发生在哪一步。

设一个教学试行的第一轮集中检查难例，发现许多文档依赖附件中的检测报告，而界面只显示正文。团队于是补充附件摘要，并增加“附件缺失”这一信息不足原因；第二轮材料变长，单例中位时间也随之上升。补证据让更多判断可执行，同时改变了每例成本，两项估计需要一起更新。第一轮为发现问题而刻意集中难例，其耗时不能直接乘生产对象数来估算总投入；扩大前还要按预计类别、长度和难度构成抽取样本，并保留新问题的反馈入口。

\section{标注方式的选择}
\label{sec:knowledge-acquisition-human-interaction}

对同一篇文档，人可以自己选出类别，也可以核对系统预填的类别；可以从头寻找证据，也可以修改已推荐的片段。任务要求未变，提供信息的动作却变了。\term{人工交互方式}（human annotation interaction）就是标注者看见什么、怎样操作、系统保存什么这一组安排。它与输出粒度可以组合：上位类别不必从零选择，细片段也不必全部由人划定。比较方案时，要一起计算阅读、判断、录入、模型推理和返工，不能把一次点击、一对偏好和一条完整答案当成等量知识。

\subsection{直接给出答案}
\label{sec:knowledge-acquisition-human-direct-answer}

如果系统只提供文档、问题和规范，由人寻找证据并产生答案，这种方式称为\term{直接作答}（direct annotation），即不依靠预填结果完成目标输出。类别集合明确时可以选项作答，具有单位与量程的测量可以填写数值，片段或区域可以用边界工具，摘要与改写则需要开放文本。选项应覆盖合理答案并保留未知与不适用；开放回答应限定事实范围、长度和可接受变体，否则后续会把不同措辞误认成不同知识。

仍用同一篇故障文档，我们可以分别要求三种输出。类别任务提交“是、否、未知、不适用”之一，未知还要附原因码；片段任务标出支持判断的最小连续文本，必要时保留多个不相邻片段；一句话摘要只复述文档内可核查事实。相同材料对应不同完成条件，不能用“都已填写”作为共同验收标准。边界任务还要规定重叠、嵌套和最小单位，例如型号名称是否包含前置品牌，使训练过程不必猜测所标范围。

理解问题、寻找证据、形成判断和录入结果，是直接作答中不同的耗时环节。长文档可能主要耗在查找证据，图像多边形可能主要耗在精细描绘，因此缩短点击步骤能节省的总时间取决于录入原本占多少。快捷键和批量录入只能减少重复操作，不能替代逐例判断。额外理由或证据确实供训练、审计或争议处理使用时，才宜设为必填；否则它既增加工作，也可能促使人事后补写并无充分依据的解释。

\subsection{提供粗粒度或局部标签}
\label{sec:knowledge-acquisition-human-partial-label}

若逐篇确定故障部件并定位全部证据超出预算，可以先询问较少信息。\term{粗粒度标签}（coarse-grained label）只缩小答案范围，例如确认“硬件问题”而暂不区分电池、电源和主板；\term{局部标签}（partial annotation）只判断对象的一部分，例如一处证据片段或一个图像关键点。两者都约束完整目标，却没有补齐完整答案。图像任务也说明了这一成本问题：对象复杂度与输出类型会改变人工工作量，成本感知选样因而需要同时估计信息和投入\scite{knowledge-acquisition-human-vijayanarasimhan2009cost}

为了让训练过程只使用人实际确认的部分，我们需要分别记录“哪些类别尚未排除”和“哪些位置已经检查”。设完整类别$y$属于细类别集合$Y$，对象$x$的粗标签依据父子类别映射返回$S(x)\subseteq Y$，确认$y\in S(x)$，但不指定集合内哪个元素为真。对位置$j=1,\ldots,\ell$，用$m_j\in\{0,1\}$标记是否检查，只有$m_j=1$时记录局部答案$a_j$。已知约束写成
\begin{equation}
  y\in S(x),
  \qquad
  \mathcal{A}_{\mathrm{local}}(x)
  =\{(j,a_j):m_j=1\}.
  \label{eq:knowledge-acquisition-human-partial-label}
\end{equation}
式~\eqref{eq:knowledge-acquisition-human-partial-label}让类别范围和已查位置各自进入训练，未确认部分仍为空缺。$m_j=0$不表示负类，$S(x)$也不能任取一个元素作为真值。保存父子类别映射、对象标识与已检查范围，正是为了防止统一数据格式时丢掉这些区别。

当人工已经确认真值属于$S(x)$时，训练可以提高模型给这个集合的总概率，而不强迫集合内各类均分概率。第~\ref{chap:knowledge-fusion}章的式~\eqref{eq:knowledge-fusion-set-loss}给出相应集合损失。它使用了“真值在集合中”这一信息，但单靠该约束仍不能分辨集合内的细类别；未经核验的候选列表也没有资格直接充当已确认约束。

图~\ref{fig:knowledge-acquisition-human-granularity}转用同一几何场景，显示局部标签究竟增加哪种信息：整图类别说明场景含目标，关键点增加对象位置，边界增加已检查对象的几何范围。最后一个面板仍有未检查区域，因此“完整”只针对已检查对象的边界，不表示全图都已核查。减少输出是否省钱，要看查找与录入各占多少时间；若寻找对象最费时，从边界改成点可能只省少量操作，若描绘最费时则可能节省较多。标签条数无法表达这项差别。

\begin{figure*}[htbp]
  \centering
  % 同一教学场景中的图像级、关键点与包围框标注；斜线区表示未检查区域。
\begin{tikzpicture}[x={\dimexpr\linewidth/182\relax},y=1mm,scale=.98,transform shape,font=\figurefont\footnotesize,text=BookInk]
  \tikzset{
    scene/.style={draw=BookRule,fill=BookPaper,line width=.8pt},
    object/.style={line width=1pt},
    selected boundary/.style={draw=FigureModel,line width=1.3pt},
    panel title/.style={font=\figurefont\bfseries\footnotesize,anchor=south},
    output/.style={draw=BookRule,fill=white,minimum width=48mm,minimum height=8mm,align=center}
  }
  \foreach \shift/\panel/\title in {0/a/图像级类别,61/b/对象关键点,122/c/对象包围框}{
    \begin{scope}[xshift=\shift mm]
      \node[panel title] at (26,49) {(\panel)\quad\title};
      \path[scene] (1,8) rectangle (51,43);
      \draw[object,draw=BookBlue,fill=FigureInput!18] (14,27) circle (6);
      \draw[object,draw=BookAmber,fill=FigureLoss!18] (25,13)--(37,13)--(31,27)--cycle;
      \draw[object,draw=BookTeal,fill=FigureModel!18] (39,30) rectangle (48,38);
    \end{scope}
  }
  \begin{scope}
    \node[output] at (26,3) {标注：场景含目标对象};
  \end{scope}
  \begin{scope}[xshift=61mm]
    \begin{scope}
      \clip (25,8) rectangle (51,43);
      \foreach \d in {-20,-15,...,70}{\draw[BookRule,line width=.55pt] (\d,8)--({\d+35},43);}
    \end{scope}
    \draw[BookMuted,dashed,line width=.8pt] (25,8) rectangle (51,43);
    \fill[FigureModel] (14,27) circle (1.5);
    \draw[white,line width=.7pt] (12.6,27)--(15.4,27) (14,25.6)--(14,28.4);
    \node[output] at (26,3) {标注：对象甲的关键点};
    \node[anchor=north east,text=BookMuted] at (49.5,41.5) {未检查};
  \end{scope}
  \begin{scope}[xshift=122mm]
    \begin{scope}
      \clip (38,8) rectangle (51,43);
      \foreach \d in {-20,-15,...,70}{\draw[BookRule,line width=.55pt] (\d,8)--({\d+35},43);}
    \end{scope}
    \draw[BookMuted,dashed,line width=.8pt] (38,8) rectangle (51,43);
    \draw[selected boundary] (7,20) rectangle (21,34);
    \draw[selected boundary] (23,11) rectangle (38,29);
    \node[output] at (26,3) {标注：已检查对象甲、乙的框};
    \node[anchor=north east,text=BookMuted] at (49.5,41.5) {未检查};
  \end{scope}
\end{tikzpicture}

  \caption{同一对象的粗粒度、局部与完整标注。（a）给整图类别，（b）给对象甲的关键点，（c）给已检查对象甲、乙的边界；三个面板共享同一教学场景。“完整”限于已检查对象的边界，斜线区域仍未知，不能作负类。几何形状不是真实识别结果，也不表示成本比例。}
  \label{fig:knowledge-acquisition-human-granularity}
\end{figure*}

\subsection{比较候选结果}
\label{sec:knowledge-acquisition-human-comparison}

要求人写出理想的故障原因概括，可能比核查两个现成回答更困难。此时可以采用\term{候选比较}（candidate comparison），即对给定结果记录偏好、排序或按维度评分。先固定评价维度$d$，再用$a\succ_d b$记录偏好$a$而非$b$，用$a\sim_d b$记录平局，用$a\parallel_d b$记录依据不足或不可比较。它们回答的是这组候选在该维度上的关系，不是候选集之外的正确答案，也不是绝对质量分数。

例如，甲完整引用诊断证据但较长，乙简短却遗漏限定条件。按“事实准确”比较可能选甲，按“简洁”比较则可能选乙。这个对照新增的是评价维度的区别，而不是给两份回答作总体排名。记录应关联问题、候选内容、展示位置、维度与作答者；只写“甲更好”，后续偏好模型便无法知道人看重的是事实还是长度。

若采用维度评分，还要让档位有可重复的含义。在事实准确的五档序数量表中，可以用“1”表示核心事实与材料冲突，“3”表示核心结论成立但遗漏重要限定，“5”表示关键事实及限定都有材料支持；“2”和“4”位于相邻锚点之间。同一回答的事实准确可以得5分，简洁性却只有2分，因为两项分数约束的属性不同。序数量表只给顺序，不保证5到3与3到1是等量质量差；没有测量模型和校准证据，也不能直接比较不同人员未对齐的打分习惯。

增加到三个以上候选后，各对结果未必组成唯一排序：可能平局，也可能出现甲优于乙、乙优于丙、丙又优于甲的循环。汇总需要另加排序规则或偏好模型，不能由一次判断直接推出总序。即使某个候选胜出，候选集也可能根本没有合格答案，所以应允许“两者都不合格”并补充问题。随机或均衡展示位置可以检查位置影响，却不能消除所有措辞、长度与知识背景差异。$J$个候选的每一无序对都比较一次，需要$J(J-1)/2$次判断；只查询部分关系可减少成本，但要保留尚未确定的顺序。

\subsection{确认或修改模型建议}
\label{sec:knowledge-acquisition-human-model-suggestion}

若模型已经找出可能的证据句，人可以先核对它，而不必从整篇文档重新搜索。\term{模型辅助标注}（model-assisted annotation）就是先由模型提出类别、边界或草稿，再由人确认、修改或拒绝。它可能节省搜索和录入，也会新增核对、定位错误和撤销建议的工作。界面应区分确认、局部编辑、完全重写、拒绝后另答与弃权；只有人明确提交未改动的建议，才能记为确认，超时或跳过都不表示认可。

这项区别要落实到记录状态。任务已发出但未开始，执行状态是“未处理”；已经开始却未提交，是“处理中”，来源动作仍为空。正常提交判断记为作答，完成检查而不判断记为弃权，两者的执行状态都可以是“成功”。成功表示操作完成，不表示内容正确或已经核验；未请求与执行失败也不能变成人主动弃权。

保存模型建议及其版本、人工操作和最终答案，才能追查错误是沿用建议还是由人另行产生。对同一文档，可以比较只显示材料的从零作答、显示多个答案的候选列表，以及默认填入一个答案的单一预填。三种界面固定任务而改变人首先看到的结果；更具体的预填可能节省更多操作，也可能让核对更依赖那个默认答案。

自动建议是否有益，需要连同具体任务检验。Schulz等在两个专家领域的篇章级任务中观察到标注速度和表现改善，在其实验中未发现值得注意的偏差；这一观察限于所研究的任务、建议模型和界面\scite{knowledge-acquisition-human-schulz2019suggestions}
某类输入若集中出现错误建议，总体平均节省仍可能掩盖该区域的系统性返工。

比较建议界面时，应固定任务、材料和最终质量检查口径，再一起记录推理成本、核对时间与修改时间。一部分随机样本可以不展示建议，让人独立作答，以检查默认答案是否改变错误模式；重要争议也应保留独立判断。人机共同答案已受该模型影响，不能又用作评价同一模型的完全独立参照。

表~\ref{tab:knowledge-acquisition-human-interactions}汇总四种方式新增的信息与留下的未知，通用来源和状态仍按章首接口保存。各行不是互斥方案：可以先确认类别建议，再从零定位证据。不过两步可能共享阅读，也可能产生额外核对，所以组合成本应重新实测，不能直接相加各自的平均时间。

\begin{table*}[htbp]
  \centering
  \small
  \begingroup
  \setlength{\tabcolsep}{3pt}
  \renewcommand{\arraystretch}{1.15}
  \begin{tabularx}{\linewidth}{@{}>{\RaggedRight\setlength{\parindent}{0pt}\arraybackslash}p{18mm}>{\RaggedRight\setlength{\parindent}{0pt}\arraybackslash}p{24mm}>{\RaggedRight\setlength{\parindent}{0pt}\arraybackslash}p{25mm}>{\RaggedRight\setlength{\parindent}{0pt}\arraybackslash}p{29mm}>{\RaggedRight\setlength{\parindent}{0pt}\arraybackslash}p{24mm}>{\RaggedRight\setlength{\parindent}{0pt}\arraybackslash}X@{}}
    \toprule
    方式 & 可见输入 & 人工操作 & 所得信息 & 仍未知 & 成本与风险 \\
    \midrule
    直接作答 & 文档、问题、规范 & 选类别、定位或撰写 & 答案及支持证据 & 未问属性、未查区域 & 阅读与生成；可能遗漏限定 \\
    粗粒度或局部判断 & 文档、类别层级或局部范围 & 选上位类、集合或局部标记 & 完整目标的约束、已查范围 & 细类别、未查位置 & 查找对象；可能把遗漏当负类 \\
    候选比较 & 问题、候选、评价维度 & 偏好、平局、不可比或评分 & 相对关系或维度评分 & 候选集外答案、未比较关系 & 比较对数增加；位置与措辞偏差、循环偏好 \\
    确认或修改建议 & 文档、建议、上下文 & 确认、编辑、重写、另答或弃权 & 原建议及版本、人工操作、最终答案 & 未独立检查的替代答案 & 生成、核对与返工；可能沿用系统性错误 \\
    \bottomrule
  \end{tabularx}
  \endgroup
  \caption{人工交互方式与所得知识。四行固定文档任务，比较可见输入、人工动作、所得信息和仍未知的部分；来源与状态字段另按章首接口保存。成本列指出需要测量的环节，不表示任何方式无条件更快或质量更高。}
  \label{tab:knowledge-acquisition-human-interactions}
\end{table*}

\section{标注样本的选择}
\label{sec:knowledge-acquisition-human-sample-selection}

现在固定故障判断任务和作答界面，候选文档仍多于预算允许的人次。我们需要决定下一批先问哪些文档，而不是继续改答案含义。\term{主动学习}（active learning）让学习系统根据已有知识选择待询问样本，即由当前学习状态参与决定新增训练材料。本章采用池式设定：第$t$轮有带标集合$L_t$、未标候选池$U_t$，以及由$L_t$训练得到的参数$\bm{\theta}^{(t)}$。打分时固定模型与可用证据，收到反馈后才更新；流式到达只改变当时可见的候选集合，不另构成一种价值准则。

怎样判断主动选择值得使用？应与随机抽样比较；已知类别、来源或时间结构重要时，还可以事先规定分层抽样。主动选样改变了各文档进入已标集合的概率，因此这些标签不能不经校正就当作总体审计样本。图~\ref{fig:knowledge-acquisition-human-query-loop}给出共同循环，后面的不确定性、代表性和预期收益主要改变打分或集合选择这一步。成功提交的可用答案进入已标集合；来源弃权只增加原因记录与已付成本，不增加确定标签。候选池也要据此更新，避免把不可答文档原样反复送回。

\begin{figure}[htbp]
  \centering
  % 池式主动选择流程；虚线路径保留来源弃权及原因，不把查询数等同于有效标签数。
\begin{tikzpicture}[x=1mm,y=1mm,font=\figurefont\footnotesize,text=BookInk,
  box/.style={draw=BookRule,fill=BookPaper,line width=.8pt,rounded corners=1.2mm,minimum width=42mm,minimum height=11mm,align=center,inner sep=2.5pt},
  decision/.style={box,draw=FigureModel,fill=FigureModel!10},
  flow/.style={knowledge arrow,draw=BookTeal,line width=1.1pt,rounded corners=1.2mm},
  secondary/.style={knowledge arrow,draw=BookAmber,line width=.85pt,dashed,rounded corners=1.2mm},
  edge label/.style={font=\figurefont\scriptsize,fill=white,inner sep=1.5pt}]
  \node[box,draw=FigureInput,fill=FigureInput!10,minimum width=31mm] (state) at (16,0) {已有标注$L_t$\\与候选池$U_t$};
  \node[box,minimum width=31mm] (criterion) at (51,0) {按当前依据\\计算查询分数};
  \node[box,minimum width=31mm] (select) at (86,0) {选择样本$x_t$\\并发出任务};
  \node[decision,minimum width=42mm] (judge) at (86,-18) {人工能否依据\\现有观测作答？};
  \node[box,draw=FigureOutput,fill=FigureOutput!7,minimum width=44mm] (usable) at (28,-36) {可用答案\\加入$L_{t+1}$并更新模型};
  \node[box,draw=FigureLoss,fill=FigureLoss!10,minimum width=45mm] (pending) at (80,-36) {来源弃权\\记录原因，$L_t$与模型不变};
  \node[decision,minimum width=54mm] (finish) at (48,-54) {更新$U_{t+1}$与累计成本后\\停止条件是否满足？};
  \node[box,draw=FigureOutput,fill=FigureOutput!7,minimum width=31mm] (end) at (95,-54) {停止并保留\\未知及其原因};
  \draw[flow] (state)--(criterion);
  \draw[flow] (criterion)--(select);
  \draw[flow] (select)--(judge);
  \draw[flow] (judge.south west)--++(0,-4)-|node[edge label,pos=.72,above] {可用}(usable.north);
  \draw[secondary] (judge.south east)--++(0,-4)-|node[edge label,pos=.72,above] {弃权}(pending.north);
  \draw[flow] (usable.south)--++(0,-4)-|(finish.north west);
  \draw[secondary] (pending.south)--++(0,-4)-|(finish.north east);
  \draw[flow] (finish)--node[edge label,above] {是}(end);
  \draw[flow] (finish.west)--node[edge label,above,pos=.25] {否}(3,-54)--(3,-12)--(51,-12)--(criterion.south);
\end{tikzpicture}

  \caption{主动选择样本的查询循环。实线路径把成功提交的可用答案加入已标集合并更新模型；虚线路径记录成功提交的来源弃权及原因，内容仍未知，已标集合和模型不变。两条路径都更新候选池与累计成本，再检查预算、覆盖或独立验证。返回箭头只表示重算依据，不保证质量提高。}
  \label{fig:knowledge-acquisition-human-query-loop}
\end{figure}

\subsection{基于不确定性的样本选择}
\label{sec:knowledge-acquisition-human-uncertainty-selection}

若模型对大部分常见故障已有稳定判断，把下一份人工答案用于它尚不能判断的文档，是一种选样思路。\term{不确定性采样}（uncertainty sampling）按当前模型的犹豫程度确定查询优先级。我们可以看一个模型如何分配类别概率，也可以看多个候选模型是否给出不同意见；两者并不等价。输入损坏、规范歧义或未知类别也可能产生高分，因此分数只提供询问依据，仍要检查可答性和覆盖，不能直接当成新增知识的保证。

\subsubsection{基于预测不确定性的样本选择}
\label{sec:knowledge-acquisition-human-predictive-uncertainty}

对同一篇候选文档，模型可能把概率集中在一个类别，也可能让几个类别难分。要把这种差别用于排序，需要一个可计算分数。固定当前模型与同一$K$类空间，令$p_c(x)=p_{\bm{\theta}^{(t)}}(y=c\mid x)$，降序概率为$p_{(1)}(x)\geq p_{(2)}(x)\geq\cdots\geq p_{(K)}(x)$，其中$K\geq2$。最低置信度（least confidence）考察第一名是否占优，类别间隔（margin）考察前两名是否接近，预测熵（predictive entropy）考察整个类别分布是否分散。把三者都写成越大越优先查询的分数：
\begin{equation}
  \begin{aligned}
    u_{\mathrm{LC}}(x)
      &=1-p_{(1)}(x),\\
    u_{\mathrm{M}}(x)
      &=1-\bigl[p_{(1)}(x)-p_{(2)}(x)\bigr],\\
    u_{\mathrm{H}}(x)
      &=-\sum_{c=1}^{K}p_c(x)\log p_c(x),
  \end{aligned}
  \label{eq:knowledge-acquisition-human-predictive-uncertainty}
\end{equation}
其中约定$0\log0=0$，本章用自然对数，熵以nat计。熵中的加权对象是该文档可能所属的类别，不是候选文档之间的变化，也不是已观测错误的平均。式~\eqref{eq:knowledge-acquisition-human-predictive-uncertainty}把三种不同的犹豫程度转成同方向分数；每轮在$U_t$中按选定分数排序，查询预算允许的候选，平分时可按预先约定随机打破。Lewis与Gale在文本分类中系统使用了顺序式不确定性采样，这里保留其查询思路并给出统一方向的多分类写法\scite{knowledge-acquisition-human-lewis1994uncertainty}

\begin{example}[三种不确定性分数可能给出不同排序]
  \label{ex:knowledge-acquisition-human-uncertainty-scores}
  为了只比较分数而不混入模型差异，设同一模型对两个候选给出三分类教学分布$\bm{p}^{A}=(0.50,0.49,0.01)$与$\bm{p}^{B}=(0.45,0.30,0.25)$，两者都已归一化。$A$的最低置信度、间隔分数和自然对数熵依次是$0.50$、$1-(0.50-0.49)=0.99$和$0.742$；$B$依次是$0.55$、$1-(0.45-0.30)=0.85$和$1.067$。$A$的前两类几乎平分，所以间隔准则选$A$；$B$的第一名更弱、整体更分散，所以最低置信度和熵选$B$。二分类中，给定最大类别概率就确定另一类概率，三种分数都随它单调变化，除平局外排序相同；多分类增加了剩余概率怎样分配的自由，等价关系不再成立。
\end{example}

若查询对象从类别文档变成逐位置标签，单个位置的犹豫还不能决定整篇文档是否值得阅读。可以对$j=1,\ldots,\ell$先算同口径局部分数$h_j$，再取总和$\sum_jh_j$、平均值$\ell^{-1}\sum_jh_j$或最大值$\max_jh_j$，以得到文档级查询分数。总和会偏向长文档，平均会稀释单个关键难点，最大值可能被一个异常位置主导。这是评价单位的改变，不是三种汇总可以随意互换。成本随长度变化时，还可比较$u(x)/\widehat c(x)$；$\widehat c(x)$必须是另行验证的人工成本估计，字符数不能直接等于耗时。

上述分数均使用模型概率。概率失准时，$p_c(x)$只表达模型的数值判断，不能解释为真实正确率。高熵可能来自类别边界，也可能来自缺页、跨语言或规范缺口；低熵也可能是模型对未知类别自信地犯错。因此应保留随机或分层覆盖，并把弃权原因送回任务设计。对同类不可答文档持续查询，只会消耗预算，不会消除缺失证据。

\subsubsection{基于模型分歧的样本选择}
\label{sec:knowledge-acquisition-human-model-disagreement}

一个模型的概率分散，不说明换一个同样合理的模型会作何判断。设已有故障标签主要覆盖常见表述，两种预测规律在这些文档上都满足当前拟合准则，却可能在少见的技术描述上作出不同预测。有限数据只约束已经观察到的情形，未必排除这些不同规律，因此我们可以保留一组候选，而不把某次训练的结果当成唯一可能解释。\term{委员会查询}（query by committee, QBC）让这组模型判断同一输入，优先询问意见不一致的文档。“与数据相容”指没有被所用数据和拟合准则排除，不意味着任意几个训练失败的模型都可以充当委员会。

怎样取得成员，也决定分歧的含义。训练子集重采样考察数据扰动，不同初始化探索优化路径，后验采样则按给定概率模型表达参数不确定性；实践中的前两种构造只是候选差异的近似，不自动等于完整后验。QBC原始工作以最大分歧选择查询，并在特定理论模型下分析信息收益，不能把那种收益无条件移到任意委员会\scite{seung1992querycommittee}

先考虑只保留最终类别的委员会。为把“意见分成几派”转成查询分数，固定同一候选$x$、同一$K$类空间及$M\geq2$个成员，第$m$个成员输出$p_m(c\mid x)$和硬预测$\widehat y_m(x)$；硬预测平票应按固定规则处理。令$v_c(x)=\sum_{m=1}^{M}\mathbb{1}\{\widehat y_m(x)=c\}$，以$v_c(x)/M$表示成员中投给类别$c$的比例。对这组票数比例取熵，得到投票熵（vote entropy），即成员选择的类别有多分散：
\begin{equation}
  u_{\mathrm{VE}}(x)
  =-\sum_{c=1}^{K}\frac{v_c(x)}{M}
    \log\frac{v_c(x)}{M}.
  \label{eq:knowledge-acquisition-human-vote-entropy}
\end{equation}
投票熵越大，成员的硬类别越分散，按此准则越应询问；它对成员投票作平均，不对真实标签错误作平均。不过两个模型即使投同一类，也可能把其概率分别设为$0.51$和$0.99$，硬票会丢掉这项差异。要使用完整分布，需要先给每个成员一个共同比较基准，即等权平均分布：
\[
  \bar p(c\mid x)
  =\frac{1}{M}\sum_{m=1}^{M}p_m(c\mid x),
\]
平均分布不是新成员的预测，也不直接证明哪个成员正确。我们再比较每个成员的概率与这个共同基准的差异，用成员概率加权对数比，然后在成员间平均：
\begin{equation}
  u_{\mathrm{KL}}(x)
  =\frac{1}{M}\sum_{m=1}^{M}
    \sum_{c=1}^{K}p_m(c\mid x)
    \log\frac{p_m(c\mid x)}{\bar p(c\mid x)}.
  \label{eq:knowledge-acquisition-human-kl-disagreement}
\end{equation}
式~\eqref{eq:knowledge-acquisition-human-kl-disagreement}是各成员到共识分布的平均Kullback--Leibler散度，也就是成员概率分配偏离共同基准的平均程度；数值越大，分布分歧越大。两种熵或散度都用自然对数，并在同一委员会、同一类别空间下比较候选。零概率项按极限取零；若某成员给正概率而平均概率为零，就要检查数值或类别映射，因为等权平均本应保留该正概率。

用三个成员判断同一批三分类候选，可以分清这两个对象。三者都在$(0.40,0.35,0.25)$附近且第一类最大时，每个分布都较分散，硬票却全选第一类，投票熵为零；三者分别给$(0.70,0.20,0.10)$、$(0.20,0.70,0.10)$和$(0.20,0.10,0.70)$时，每类得到一票，投票熵达到这个三成员三分类设定的最大值。若都给$(0.90,0.06,0.04)$，单模型熵与模型分歧都较低。这些教学概率说明，类别概率分散和成员意见不一是不同判断，不能只凭其中一个反推另一个。

保留这种分歧信息，需要训练或保存$M$个模型、对每个候选预测$M$次，并在新标签到来后刷新成员。成员还可能因相同数据或架构共享偏差，一致地犯错；初始化差异也只覆盖部分优化差异，不是完整参数后验。Settles与Craven在序列标注中比较了委员会、预期梯度长度等策略，具体排序依赖任务、模型和成本口径，因此不能从该比较指定一个普遍最优准则\scite{knowledge-acquisition-human-settles2008analysis}

\subsection{基于代表性与多样性的样本选择}
\label{sec:knowledge-acquisition-human-representative-diverse-selection}

假设候选池有大量近似的常见投诉，还有少量不同部件的报告。只挑不确定性最高的文档，可能一次问到许多几乎重复的内容。我们还要看一个标签能覆盖哪些相似情形，以及批内是否重复。\term{代表性}（representativeness）考察候选是否处在数据的重要区域，\term{多样性}（diversity）考察所选一批是否包含不同情形。判断相似先要把文档映射成表示$\bm{\phi}(x)$，再定义距离或相似度；表示若主要反映长度或模板，算法获得的只是那些维度的覆盖，不一定是故障语义的覆盖。

先用候选池估计哪些区域常见。给定固定的非负相似度$K(\bm{\phi}(x),\bm{\phi}(z))$，数值越大表示两篇文档越相似；这里$K(\cdot,\cdot)$是相似度函数，不是前面的类别数。对非空$U_t$，把$x$与池中每个$z$的相似度平均，得到局部密度分数：
\begin{equation}
  \rho(x)=\frac{1}{|U_t|}
  \sum_{z\in U_t}K\bigl(\bm{\phi}(x),\bm{\phi}(z)\bigr).
  \label{eq:knowledge-acquisition-human-density}
\end{equation}
高$\rho(x)$表示与$x$相似的候选较多，因此优先它容易覆盖常见区域，也可能遗漏稀有类别。批量选择还需回答另一问题：已有标签和本批标签合在一起，池中每篇文档能找到多相似的代表？令$X_t=\{x:(x,y)\in L_t\}$为已标输入，允许批次是$U_t$中所有大小为$b$的子集，$1\leq b\leq|U_t|$。对每个$z$只取它与$X_t\cup B$中最相似样本的相似度，再对池内文档求和；最大化这个总覆盖，得到
\begin{equation}
  B^\ast\in\argmax_{\substack{B\subseteq U_t\\|B|=b}}
  \sum_{z\in U_t}
  \max_{x\in X_t\cup B}
  K\bigl(\bm{\phi}(x),\bm{\phi}(z)\bigr).
  \label{eq:knowledge-acquisition-human-coverage}
\end{equation}
内层最大值让每个$z$只由最相似代表贡献一次覆盖，避免把多个代表与同一文档的相似度重复相加；外层求和衡量池内总体覆盖，而非最坏文档的覆盖保证。在教学设定中，某文档当前最大相似度为$0.8$，新样本与它相似度为$0.9$，该项增益是$\max(0.8,0.9)-0.8=0.1$；已有覆盖若为$0.95$，增益就变成$\max(0.95,0.9)-0.95=0$。随着已选集合增大，同一新样本的覆盖增益不会增加，这种边际收益递减称为\term{次模性}（submodularity）。所以每选一个样本都应重算它对当前$X_t\cup B$的新增覆盖；直接取密度最高的$b$个，仍可能集中在同一簇。

式~\eqref{eq:knowledge-acquisition-human-coverage}称为facility-location型覆盖目标，它用已选代表与其余文档的相似度评价一组样本。目标定义了怎样比较批次，还没有给出怎样找到最优批次。这类次模目标已用于数据子集选择和主动学习\scite{knowledge-acquisition-human-wei2015submodularity}

为计算每次加入样本的增益，我们把任意代表集合$S$的总覆盖记为$F(S)$，再减去固定已标集合的覆盖，得到批次新增覆盖$G(B)$：
\[
  F(S)=\sum_{z\in U_t}\max_{x\in S}
  K\bigl(\bm{\phi}(x),\bm{\phi}(z)\bigr),
  \qquad
  G(B)=F(X_t\cup B)-F(X_t).
\]
约定空代表集合的最大相似度为0，就有$G(\varnothing)=0$，即尚未增加样本时增益为零；这称为归一化。增加代表不会降低原有最大相似度，因此$G$单调，再结合边际收益递减得到归一化、单调的次模函数。一般问题包含最大覆盖问题，精确求解是NP-hard，即不存在已知的通用高效精确算法。实践中可从$B=\varnothing$出发，每轮在未选候选中加入使$G(B\cup\{x\})-G(B)$最大的一个，直至选满$b$个。相似度固定且非负、只用批量大小的基数约束时，标准贪心至少达到最优覆盖增益的$1-1/e$；保证对象是相似度覆盖，不是分类准确率\scite{knowledge-acquisition-human-nemhauser1978submodular}

图~\ref{fig:knowledge-acquisition-human-selection-criteria}固定同一二维教学点集，让我们比较三种依据改变了哪些选择。靠近给定边界的点体现预测不确定性，高密度区域内的候选体现代表性，不同区域的一批点体现多样性。右上孤立点既可能是稀有类别，也可能是异常，距离本身不能分辨。池式查询只允许选择实际候选；聚类的几何中心若不对应真实文档，就只能帮助寻找附近样本，不能直接交给人作答。

\begin{figure*}[htbp]
  \centering
  % 三种选样准则共享完全相同的教学点集、坐标范围和孤立点。
\begin{tikzpicture}[x={\dimexpr\linewidth/174\relax},y=1mm,font=\figurefont\footnotesize,text=BookInk]
  \tikzset{
    candidate/.style={circle,draw=BookMuted,fill=white,line width=.7pt,minimum size=3.2mm,inner sep=0pt},
    chosen/.style={circle,draw=FigureModel,fill=FigureModel,line width=1pt,minimum size=3.8mm,inner sep=0pt},
    panel title/.style={font=\figurefont\bfseries\footnotesize,anchor=south},
    criterion/.style={font=\figurefont\scriptsize,text=BookMuted,align=center},
    leader/.style={draw=BookMuted,line width=.65pt}
  }
  \def\candidatepoints{5/9,8/13,10/7,12/17,15/11,20/27,24/10,24/31,27/15,29/29,30/8,33/19,36/13,45/34}
  \foreach \shift/\panel/\title in {0/a/预测不确定性,61/b/代表性,122/c/批内多样性}{
    \begin{scope}[xshift=\shift mm]
      \node[panel title] at (26,46) {(\panel)\quad\title};
      \draw[BookRule,line width=.7pt] (2,5) rectangle (50,40);
      \foreach \x/\y in \candidatepoints{\node[candidate] at (\x,\y) {};}
      \node[criterion] at (26,1.5) {\strut 候选为空心，选中样本为实心};
    \end{scope}
  }
  \begin{scope}
    \draw[FigureLoss,dashed,line width=1pt] (22,5)--(22,40);
    \node[chosen] at (20,27) {};
    \node[chosen] at (24,10) {};
    \node[anchor=south west,text=FigureLoss] at (22.5,35.5) {给定边界};
    \node[anchor=north east,text=BookMuted] at (47,33) {孤立点};
  \end{scope}
  \begin{scope}[xshift=61mm]
    \draw[FigureInput!35,line width=.7pt] (10,11) ellipse (8 and 6);
    \draw[FigureInput!35,line width=.7pt] (31,13) ellipse (8 and 7);
    \draw[FigureInput!35,line width=.7pt] (25,30) circle (7);
    \node[chosen] at (8,13) {};
    \node[chosen] at (27,15) {};
    \node[chosen] at (24,31) {};
    \node[anchor=north east,text=BookMuted] at (47,33) {孤立点};
  \end{scope}
  \begin{scope}[xshift=122mm]
    \node[chosen] at (5,9) {};
    \node[chosen] at (29,29) {};
    \node[chosen] at (36,13) {};
    \draw[FigureLoss,line width=.8pt] (43,32)--(47,36) (43,36)--(47,32);
    \draw[leader] (38.5,35.5)--(43.2,34.3);
    \node[anchor=north east,text=BookMuted,align=right] at (38,39) {孤立点\\未必有用};
  \end{scope}
\end{tikzpicture}

  \caption{不确定性、代表性与多样性选择的差别。三面板共享教学点集与坐标范围，空心圆为候选，实心圆为池内选中样本。（a）按距给定虚线边界的接近程度选择；（b）从三个高密度区域各选一个实际候选，浅轮廓只标区域；（c）选覆盖不同区域的一批，并用孤立点说明距离远未必有用。坐标不是真实类别或性能证据。}
  \label{fig:knowledge-acquisition-human-selection-criteria}
\end{figure*}

两种依据可以组合，但先筛选什么会改变允许集合。先按可答性和密度排除损坏输入及极端孤立点，再按不确定性排序，是以这些条件决定资格；先保留高不确定候选，再最大化批内覆盖，则以不确定性决定资格。后一种仍要独立检查可答性。若直接线性加权，必须处理分数尺度并验证权重；给出两个权重本身不支持普遍最优，也不能证明被排除的孤立点没有价值。

覆盖带来的计算应分两步核算。在固定表示和单次相似度计算成本的口径下，显式构造$|U_t|\times|U_t|$矩阵需$O(|U_t|^2)$时间与存储；朴素贪心的每轮还要扫描剩余候选，并对全部$z\in U_t$算边际增益，最坏为$O(b|U_t|^2)$。缓存当前最大相似度、用懒惰贪心避免重复评估，或先预筛候选，可以减少实际计算，但不改变目标的组合优化性质。近似近邻或预筛也可能漏掉稀有区域，需要检查其覆盖代价。

\subsection{基于预期收益的样本选择}
\label{sec:knowledge-acquisition-human-expected-utility-selection}

假设两篇文档同样不确定，却只有一篇的标签会影响大量待处理报告，单看当前概率无法分辨它们的后续价值。预期收益方法因而问：得到标签并更新模型后，目标误差或参数预计改变多少？标签现在未知，不能只挑最有利的答案来计算；必须逐个假设可能标签，再按查询前的模型概率取条件期望。预期误差降低比较目标分布上的损失，预期模型变化比较参数或梯度幅度。前者直接面对任务误差却计算昂贵，后者省去部分计算，却只是一种变化代理。以下先按返回可用且标签准确的理想条件计算，弃权、噪声和人工费用仍要另行纳入实际决策。

\subsubsection{基于预期误差降低的样本选择}
\label{sec:knowledge-acquisition-human-expected-error-reduction}

要比较“标这篇之后是否少出错”，先需要在计算中临时训练各假设标签对应的模型。这里回到分类任务，$K$重新表示类别数，不再是前节的相似度函数。固定训练过程$\mathcal{A}$，包括模型设定与求解口径，当前参数为$\bm{\theta}^{(t)}=\mathcal{A}(L_t)$。对候选$x\in U_t$和每个允许类别$c$，仅在临时数据中加入$(x,c)$，用同一过程更新：
\begin{equation}
  \bm{\theta}^{x,c}
  =\mathcal{A}\bigl(L_t\cup\{(x,c)\}\bigr).
  \label{eq:knowledge-acquisition-human-hypothetical-update}
\end{equation}
这些$\bm{\theta}^{x,c}$分别回答“若标签是$c$，会得到哪一个模型”，不会在真实带标集合中同时加入所有假设。接下来需要在共同材料上比较它们的错误，而非比较各自训练损失。取与目标使用场景相符、包含$m>0$个样本的评估集$E=\{(\bm{z}_j,r_j)\}_{j=1}^{m}$，其中$\bm{z}_j$是输入、$r_j$是参考标签，$f_{\bm{\theta}}$是该参数下的预测函数。损失$\ell$给预测与参考的偏差计分，越小越好；对同一$E$求平均得到经验风险：
\begin{equation}
  \widehat R_E(\bm{\theta})
  =\frac{1}{m}\sum_{j=1}^{m}
  \ell\bigl(f_{\bm{\theta}}(\bm{z}_j),r_j\bigr).
  \label{eq:knowledge-acquisition-human-evaluation-risk}
\end{equation}
式~\eqref{eq:knowledge-acquisition-human-evaluation-risk}给每个模型一个同口径的平均损失；取零一损失时是错误率，取对数损失时则评价概率预测，不能混用两种尺度。若$E$没有真实标签，不能直接代入已知$r_j$，可以像Roy与McCallum的池式方法那样，用当前或临时模型的预测分布估计对数损失或零一损失。这样又增加了假设：模型概率与候选池足以近似真实评估分布\scite{knowledge-acquisition-human-roy2001error}

各分支风险仍不能单独决定查询，因为人工只会返回一个真实标签。我们用查询前的$p_{\bm{\theta}^{(t)}}(c\mid x)$估计各互斥类别的可能性，对分支风险加权平均，再从当前风险减去它。\term{预期误差降低}（expected error reduction, EER）就是这项预计减少的平均损失，其分支汇总与查询分数为
\begin{align}
  \overline R_E(x)
    &=\sum_{c=1}^{K}
      p_{\bm{\theta}^{(t)}}(c\mid x)
      \widehat R_E(\bm{\theta}^{x,c}),
  \label{eq:knowledge-acquisition-human-expected-future-risk}\\
  u_{\mathrm{EER}}(x)
    &=\widehat R_E(\bm{\theta}^{(t)})
      -\overline R_E(x).
  \label{eq:knowledge-acquisition-human-expected-error-reduction}
\end{align}
本轮在$U_t$内最大化$u_{\mathrm{EER}}$，等价于最小化$\overline R_E$，因为当前风险对所有候选相同。保留这个基线，可区分“预计更新后风险仍很高”和“相对当前确有下降”；分数也可能为负，表示预计损失增加。权重只是查询前模型对未知标签的估计，不是已观测频率，所以计算所得下降不保证实际查询后也下降。

\begin{example}[可复算的预期误差降低]
  \label{ex:knowledge-acquisition-human-expected-error}
  设同一评估集和损失下，当前风险为$0.30$，以下均为二分类教学数值。$x_A$的标签概率是$0.6$和$0.4$，对应临时更新后的风险是$0.18$和$0.27$，所以$\overline R_E(x_A)=0.6\times0.18+0.4\times0.27=0.216$，降低量为$0.30-0.216=0.084$。$x_B$两类各占$0.5$，分支风险是$0.12$和$0.35$，得到$\overline R_E(x_B)=0.235$、降低量$0.065$。按EER应选$x_A$，虽然$x_B$有一个更有利的分支。决定排序的是所有未知标签分支的加权结果，不是最有利标签下的最低风险。
\end{example}

图~\ref{fig:knowledge-acquisition-human-expected-utility}把$x_A$的两种互斥假设分开画出：各自临时更新、各自估计风险，最后才加权。图中两份计算不意味着向人索取两份互相冲突的真值。$N$个候选、$K$类的朴素计算需要$NK$次临时训练与评估；预筛候选、抽样标签、增量更新或近似重训可以降低开销，却也改变分数，仍需用实际查询后的独立任务效果检验。

\begin{figure*}[htbp]
  \centering
  % 二分类候选的互斥标签分支，以及不做完整重训的期望梯度长度替代路径。
\begin{tikzpicture}[x={\dimexpr\linewidth/174\relax},y=1mm,font=\figurefont\footnotesize,text=BookInk,
  box/.style={draw=BookRule,fill=BookPaper,line width=.8pt,rounded corners=1.2mm,minimum height=11mm,align=center,inner sep=2.5pt},
  flow/.style={knowledge arrow,draw=BookTeal,line width=1.1pt,rounded corners=1.2mm},
  alternative/.style={knowledge arrow,draw=BookAmber,line width=.9pt,dashed,rounded corners=1.2mm},
  edge label/.style={font=\figurefont\scriptsize,fill=white,inner sep=1.3pt}]
  \node[box,draw=FigureInput,fill=FigureInput!10,minimum width=25mm] (candidate) at (15,25) {候选样本$x$};
  \node[box,minimum width=34mm] (update1) at (58,37) {假设$y=1$\\临时更新$\bm{\theta}^{x,1}$};
  \node[box,minimum width=28mm] (risk1) at (103,37) {估计风险\\$0.18$};
  \node[box,minimum width=34mm] (update0) at (58,14) {假设$y=0$\\临时更新$\bm{\theta}^{x,0}$};
  \node[box,minimum width=28mm] (risk0) at (103,14) {估计风险\\$0.27$};
  \node[box,draw=FigureOutput,fill=FigureOutput!7,minimum width=39mm] (aggregate) at (151,25) {期望风险$0.216$\\降低量$0.084$};
  \node[font=\figurefont\scriptsize,text=BookMuted] (current) at (151,43) {当前风险$0.30$};
  \draw[flow] (candidate.east)--node[edge label,above,pos=.46] {$p(y=1\mid x)=0.6$}(update1.west);
  \draw[flow] (candidate.east)--node[edge label,below,pos=.46] {$p(y=0\mid x)=0.4$}(update0.west);
  \draw[flow] (update1)--(risk1);
  \draw[flow] (update0)--(risk0);
  \draw[flow] (risk1.east)--(aggregate.west);
  \draw[flow] (risk0.east)--(aggregate.west);
  \draw[flow] (current)--node[edge label,right] {比较}(aggregate);
  \node[box,draw=FigureLoss,fill=FigureLoss!9,minimum width=45mm] (gradients) at (66,-8) {分别计算$\lVert\bm{g}_1(x)\rVert_2$\\与$\lVert\bm{g}_0(x)\rVert_2$};
  \node[box,draw=FigureLoss,fill=FigureLoss!9,minimum width=47mm] (egl) at (128,-8) {$0.6\lVert\bm{g}_1\rVert_2+0.4\lVert\bm{g}_0\rVert_2$};
  \draw[alternative] (candidate.south)|-node[edge label,below,pos=.4] {不作完整重训}(gradients.west);
  \draw[alternative] (gradients.east)--(egl.west);
\end{tikzpicture}

  \caption{标签未知时的预期查询收益。实线使用例~\ref{ex:knowledge-acquisition-human-expected-error}的二分类教学数值：互斥标签各自临时更新与评估，按$0.6$和$0.4$汇总成风险$0.216$，比当前$0.30$减少$0.084$。下方虚线是下一方法的替代计算，先取各假设标签的梯度范数，再按概率平均，不先平均可能相反的梯度方向。}
  \label{fig:knowledge-acquisition-human-expected-utility}
\end{figure*}

有标签$E$参与反复选样，就属于策略使用过的验证证据，最终验收应另留独立材料，避免把对$E$的适配误报成新数据收益。无标签池的风险近似则依赖模型概率；池与目标分布不同，估计可能系统失真。单点计算还固定了当前$L_t$，两篇相似文档分别有用，不表示一起加入时收益相加，因为第一篇可能已消除第二篇的大部分信息缺口。

\subsubsection{基于预期模型变化的样本选择}
\label{sec:knowledge-acquisition-human-expected-model-change}

若为每个候选和标签重训、再遍历评估集的成本过高，可以先估计这个标签会让当前参数移动多少。\term{预期模型变化}（expected model change）用预计更新幅度作为查询依据，它不是直接估计未来错误率。为在不完整重训的条件下取得可计算代理，固定当前参数与可微的单样本损失$\ell(f_{\bm{\theta}}(x),c)$，对每个假设类别$c$求梯度：
\begin{equation}
  \bm{g}_c(x)
  =\left.
    \nabla_{\bm{\theta}}
    \ell(f_{\bm{\theta}}(x),c)
    \right|_{\bm{\theta}=\bm{\theta}^{(t)}}.
  \label{eq:knowledge-acquisition-human-hypothetical-gradient}
\end{equation}
梯度$\bm{g}_c(x)$给出这个假设标签在当前参数处的损失上升方向，普通梯度下降沿其负方向更新。若用学习率$\eta>0$作一步普通梯度更新，以$\bm{\theta}^{x,c}\approx\bm{\theta}^{(t)}-\eta\bm{g}_c(x)$近似前面的重训结果，位移长度就由$\eta\lVert\bm{g}_c(x)\rVert_2$近似。由于标签仍未知，我们先取每个分支的长度，再按查询前类别概率平均，得到\term{预期梯度长度}（expected gradient length, EGL），即预计一次单样本梯度有多大：
\begin{equation}
  u_{\mathrm{EGL}}(x)
  =\sum_{c=1}^{K}
    p_{\bm{\theta}^{(t)}}(c\mid x)
    \lVert\bm{g}_c(x)\rVert_2.
  \label{eq:knowledge-acquisition-human-expected-gradient-length}
\end{equation}
在同一参数化、同一当前状态和共同学习率下，$\eta u_{\mathrm{EGL}}(x)$近似预期一步参数位移，所以可在候选池中优先询问高EGL样本。它回答更新幅度而非更新是否有利。Settles与Craven将这一思路用于序列标注，并讨论未知标签的期望计算\scite{knowledge-acquisition-human-settles2008analysis}

为什么必须先取范数？若先平均梯度再取范数，$\lVert\sum_c p(c\mid x)\bm{g}_c(x)\rVert_2$会把相反方向抵消，衡量的是平均方向长度，不是各分支的平均移动长度。沿用\BookSectionRef{sec:classification-logistic-learning}{卷4《模型：经典模型》第3章“分类”第3.4.1.2节“参数学习”}中的二分类损失，把文档特征末尾加常数1以容纳偏置，记为增广特征$\widetilde{\bm{x}}$；单样本梯度是预测概率与假设标签之差乘这个特征。当$p(y=1\mid x)=0.5$时，两分支梯度是$-0.5\widetilde{\bm{x}}$和$0.5\widetilde{\bm{x}}$。平均方向为零，EGL却是$0.5\lVert\widetilde{\bm{x}}\rVert_2$，因为无论取得哪个标签都会发生非零更新。

再看幅度由什么造成。令$p=p(y=1\mid x)$，Logistic两分支的梯度范数分别为$(1-p)\lVert\widetilde{\bm{x}}\rVert_2$和$p\lVert\widetilde{\bm{x}}\rVert_2$，按$p$与$1-p$加权就得到$u_{\mathrm{EGL}}(x)=2p(1-p)\lVert\widetilde{\bm{x}}\rVert_2$。教学设定中，$p=0.5$、特征范数1的边界样本得$0.5$；$p=0.9$、异常特征范数10的高置信样本反而得$1.8$。后者高分主要来自尺度，未必更有价值。标准化、梯度裁剪或只考察部分参数都会改变所比较的幅度，应固定参数化和优化状态后再解释分数。

EGL不必对每个分支完整重训和扫描评估集，但仍要对每个候选、每个可能标签反向传播。动量、自适应预条件或二阶方法还会利用优化器状态改变实际步长，所以欧氏梯度范数只能是代理，不能直接等同于它们的参数位移。最终应检查查询后的独立任务损失或其他预定指标，而不是把参数变化大本身当成成功。

\section{标注者的选择}
\label{sec:knowledge-acquisition-human-annotator-selection}

一篇写着“无法开机”的中文文档，和一篇需要核对专业检测报告的外文文档，未必适合交给同一个人。领域知识、语言、工具熟练度、证据访问权限和当前负载，都会影响逐篇作答。我们先用任务要求确定谁有资格承接，再用历史结果补充能力证据，最后把样本与人员匹配。这里决定谁做什么；投入多少人次与怎样融合多人答案，分别由第~\ref{sec:knowledge-acquisition-human-budget}节和第~\ref{chap:knowledge-fusion}章回答。

\subsection{基于专业能力的标注者选择}
\label{sec:knowledge-acquisition-human-expertise-selection}

对需要核对检测报告的文档，只有“某领域专家”这一身份还不足以确定可否承接。任务需要检查四类具体能力：是否能识别目标对象、理解所用概念、查阅允许证据，以及按规范录入。领域资历支持准入，却不证明熟悉当前类别边界；受训普通人员也可能正确执行规则明确的常见判断，但缺少专业知识时应暴露缺口，而不是猜测结论。

培训例题与试标可以逐项验证这些要求。明确的普通故障描述可由理解规范的受训人员处理；需要判读电池热失控机理并核查报告的文档，则要求相应知识和访问权限。补充术语表与证据摘要后，普通人员若已能依规则判断，就不必仅凭身份交给专家。相反，任务若要求对冲突医学证据作专业裁决，短期操作培训不能替代相应专业判断。后一例改变了任务领域，只用于说明准入依据仍须由具体判断要求决定。

是否可以使用非专家，要在具体任务及答案聚合条件下验证。Snow等对若干自然语言处理任务比较非专家标注与专家标准，展示了多份非专家答案在其研究任务中的可用性；这不证明任意复杂任务都能替代专家，也不使多数一致天然正确\scite{knowledge-acquisition-human-snow2008nonexpert}
若目标转为偏好测量，所需人群也随之改变：应询问预定服务对象，记录他们的判断分布，而不是让专家替目标人群表达偏好。

人员准入也消耗资源，培训、资格试题和更新后的再确认都应计入准备成本。旧资格只支持过去任务版本；类别边界、工具或证据来源实质变化时，应重查受影响的例题。界面还应允许人报告“缺少知识”，这样才能触发转交；强制作答会把已知能力缺口隐藏成一个看似确定的答案。

\subsection{基于历史表现的标注者选择}
\label{sec:knowledge-acquisition-human-history-selection}

假设甲主要做普通故障描述，乙主要做技术报告，直接比较总体符合率可能把任务难度当成人员差异。历史记录除通用来源和状态外，还要关联任务类型、独立参考或后续裁决，以及处理时长。选择人员时，应优先比较重叠任务，或按可解释的类别和难度分组；只做简单样本的人总体比例较高，并不能证明他更适合困难样本。

在同一任务组$g$和规范版本下，设标注者$a$完成$n_{a,g}>0$项，其中$k_{a,g}$项与独立参考一致。用$\widehat q_{a,g}=k_{a,g}/n_{a,g}$描述这组已完成任务的参考符合比例，而不把它当成所有未来任务的正确率。$9/10$与$90/100$点估计相同，证据量却不同；还应并列看弃权、超时和版本，不能让一个综合分数遮住某组缺少证据的事实。可靠性估计及统计不确定性由第~\ref{chap:knowledge-quality-assessment-correction}章展开。

新加入者没有足够历史，长期无记录或任务已变更的人员也不能直接沿用旧结论。可先分配少量有已知参考、难度覆盖真实任务构成的校准材料，并在积累证据前限制高风险任务量，避免凭几次结果永久定级。参考也须检查：多数人可能共享规范误解，过时或只含常见类别的参考集则不能支持困难区域的能力判断。

\subsection{基于样本与标注者匹配的任务分配}
\label{sec:knowledge-acquisition-human-matching}

若本批包含不同语言和专业方向的文档，拥有合格人员还不等于每篇都有可用承接人。需要把逐样本资格与本轮工时写成可检查约束。对本轮要求分配的样本$i$和人员$a$，令$e_{i,a}\in\{0,1\}$表示语言、领域、权限和任务资格是否满足，$z_{i,a}\in\{0,1\}$表示是否实际分配。预计工时$t_{i,a}$与可用总工时$H_a$采用同一单位，例如小时；基本约束为
\begin{equation}
  z_{i,a}\leq e_{i,a},
  \qquad
  \sum_a z_{i,a}\geq1,
  \qquad
  \sum_i t_{i,a}z_{i,a}\leq H_a,
  \label{eq:knowledge-acquisition-human-assignment}
\end{equation}
第一项排除无资格分配，第二项使每个本轮样本至少有一名承接人，第三项避免预计工时超过个人容量。它们检验可行性，并不优化人员排名或保证得到答案；没有合格承接人或容量不足时，这组约束可能无解。允许弃权时，分配至少一次仍不等于取得可用标签，因此还要指定转交人员及等待上限，到期未解决则保留未知与原因。

图~\ref{fig:knowledge-acquisition-human-matching}用混合领域文档显示这项逐例匹配：普通中文材料交受训通用人员，医疗证据交医疗专家，英文合同交双语法律人员。这里更换领域是为了显出语言与专业资格的对应，不继承故障任务的类别规则。细线仅表示可承接，粗线才是实际分配，虚线箭头是通用人员遇专业难例后弃权转交；图中“余1项”等名额是按该教学任务预计耗时折算的负载提示。总体历史均分既不能表达这些资格，也不能表达专家只剩一个名额。

\begin{figure}[htbp]
  \centering
  % 样本需求与标注者能力的候选关系、实际分配和弃权转交。
\begin{tikzpicture}[x=1mm,y=1mm,font=\figurefont\footnotesize,text=BookInk,
  sample/.style={draw=FigureInput,fill=FigureInput!10,line width=.9pt,rounded corners=1.2mm,minimum width=38mm,minimum height=12mm,align=center,inner sep=2.5pt},
  worker/.style={draw=FigureModel,fill=FigureModel!10,line width=.9pt,rounded corners=1.2mm,minimum width=39mm,minimum height=12mm,align=center,inner sep=2.5pt},
  candidate/.style={draw=BookMuted!65,line width=.65pt},
  assigned/.style={draw=BookTeal,line width=1.5pt},
  transfer/.style={knowledge arrow,draw=BookAmber,line width=.9pt,dashed,rounded corners=1.2mm}]
  \node[font=\figurefont\bfseries\small] at (20,51) {样本需求};
  \node[font=\figurefont\bfseries\small] at (89,51) {标注者能力与负载};
  \node[sample] (s1) at (20,36) {普通中文文档\\低难度};
  \node[sample] (s2) at (20,16) {医疗证据片段\\需查阅术语};
  \node[sample] (s3) at (20,-4) {英文合同条款\\高难度};
  \node[worker] (w1) at (89,36) {受训通用标注者\\中文，余6项};
  \node[worker] (w2) at (89,16) {医疗领域专家\\中文，余1项};
  \node[worker] (w3) at (89,-4) {双语法律人员\\中英，余2项};
  \draw[candidate] (s1.east)--(w1.west);
  \draw[candidate] (s2.east)--(w1.west);
  \draw[candidate] (s2.east)--(w2.west);
  \draw[candidate] (s3.east)--(w3.west);
  \draw[assigned] (s1.east)--(w1.west);
  \draw[assigned] (s2.east)--(w2.west);
  \draw[assigned] (s3.east)--(w3.west);
  \draw[transfer] (w1.east)--(112,36)--node[left,align=right] {难例弃权\\转交}(112,16)--(w2.east);
  \draw[candidate] (8,-18)--(27,-18);
  \node[anchor=west] at (29,-18) {可承接};
  \draw[assigned] (58,-18)--(77,-18);
  \node[anchor=west] at (79,-18) {本轮分配};
  \draw[transfer] (8,-27)--(27,-27);
  \node[anchor=west] at (29,-27) {来源弃权后的转交};
\end{tikzpicture}

  \caption{样本需求与标注者能力的匹配。细实线表示具备承接条件，粗实线表示本轮分配，带箭头虚线表示来源弃权后的转交。领域、语言、难度和剩余名额均为教学假设；名额需按预计耗时换算为工时，连线数量不表示能力排名。}
  \label{fig:knowledge-acquisition-human-matching}
\end{figure}

常见样本构成稳定时，可以预设通用与专业队列；任务变化后，则应结合新答案、弃权和负载调整分工。处理较少的人员也需适量校准任务，以更新能力证据。若只把容易任务交给当前高分人员，并持续推迟难例，困难区域会长期缺少标签，历史比例也会越来越受分配方式影响。记录每次调整的依据和版本，才能区分人员能力变化与任务构成变化。

\section{标注预算的分配}
\label{sec:knowledge-acquisition-human-budget}

现在假设我们能做几类行动：询问新文档、给已有文档追加独立判断、调用专家、先粗标再细化，或启动下一轮模型更新。剩余预算应给哪类行动、各做多少，不能只靠某个样本分数回答。为比较完整配置，设$z_k$为行动$k$的非负整数执行次数，$c_k$为单次费用，$F(\bm{z})$为该配置启动的固定费用；费用均用同一种货币或约定费用单位。资源$r$可以是专家、通用人员或受限计算队列，$t_{k,r}$与容量$H_r$按该资源占用小时计，人工资源对应人工工时，不把计算占用与人工工时相加。以任务证据估计的总价值$\widehat V(\bm{z})$越大越好，规划形式为
\begin{equation}
  \begin{aligned}
    \max_{\bm{z}}\quad
      &\widehat V(\bm{z})\\
    \text{s.t.}\quad
      &F(\bm{z})+\sum_k c_k z_k\leq B,\\
      &\sum_k t_{k,r}z_k\leq H_r
      \quad\text{对每类受限资源$r$}.
  \end{aligned}
  \label{eq:knowledge-acquisition-human-budget-allocation}
\end{equation}
式~\eqref{eq:knowledge-acquisition-human-budget-allocation}在费用预算$B$和各资源总工时内选择价值尽可能高的配置，$\widehat V$可描述预期信息或预定任务收益，但要固定其口径。固定费用只对启动的配置计算，人次则由$z_k$及每项行动的人员组成决定，不能与费用或工时混为一项。总工时也不是交付期限：两个独立任务各需1小时，两人并行可在1小时内完成，却消耗2工时；普通标注与专家复核各需1小时但必须串行，至少2小时才能交付。期限还需另查并行安排、依赖和等待。

一次培训可能使许多后续标注可执行，两份相似答案则可能重复提供信息，所以行动价值不必逐项相加。固定启动费用也会改变批量是否值得开展。单位费用价值可帮助局部比较，却不保证完整组合最优；以下分别考察覆盖与重复、普通与专家、粒度与轮次的取舍。

\subsection{新增标注与重复标注的预算分配}
\label{sec:knowledge-acquisition-human-new-versus-repeat}

假设已有一些争议故障文档，而另一些部件尚无标签。把下一人次用于新文档，补的是对象覆盖；交给另一人检查旧文档，补的是同一对象的独立证据。设样本$i$有$r_i$次判断，第$j$次成本为$c_{i,j}$，逐例费用为$\sum_{j=1}^{r_i}c_{i,j}$，总费用还要加准备、培训与协调等固定项。相同总人次下，更多$r_i=1$扩大覆盖，提高少数$r_i$可观察分歧，并在适当条件下降低独立随机错误。取舍仍由遗漏类别、歧义、人员能力和下游损失决定，不由标签条数决定。

\Needspace{65mm}
第二份答案何时增加证据，取决于错误是否独立。多人共享错误规范，会重复同一个偏差；后答者看到前答者结果，也不再构成独立判断。Sheng等对噪声标注的研究表明，重复在一些质量与成本条件下可改善标签和模型，但并非总有收益，有选择地追加通常比无差别重复更值得考虑\scite{knowledge-acquisition-human-sheng2008another}
因此重复次数应成为预算决策，而不是预设“每例三人”的普遍规则。

固定次数容易执行，却对清晰与争议样本投入相同人次。自适应安排可先取得一份或两份独立答案，再按预设的分歧、证据缺失或高风险触发条件追加，并规定最高人次和停止去向。追加前要检查原因：缺失证据需要补材料，规范缺陷需要修正规则并重查，不能靠无限投票消除。多人答案的融合属于第~\ref{chap:knowledge-fusion}章，独立审计属于第~\ref{chap:knowledge-quality-assessment-correction}章，两项资源都应在当前预算中另行预留。

\subsection{普通标注与专家标注的预算分配}
\label{sec:knowledge-acquisition-human-generalist-versus-expert}

对规则清楚的故障报告和少量专业难例，可以全程请专家，也可以普通人员先处理、困难时转交。要比较费用，先固定$n$个样本的同一完成要求，并用统一费用单位记普通单价$c_{\mathrm{O}}$、专家单价$c_{\mathrm{E}}$和每次转交成本$c_{\mathrm{X}}$。全程专家的固定启动费为$F_{\mathrm{E}}$，总费用是$C_{\mathrm{all}}=F_{\mathrm{E}}+nc_{\mathrm{E}}$。普通人员先处理全部样本，再把比例$q$转专家时，普通费用并不会退回；加上固定路由费用$F_{\mathrm{R}}$，得到
\begin{equation}
  C_{\mathrm{route}}
  =F_{\mathrm{R}}+nc_{\mathrm{O}}
   +qn(c_{\mathrm{E}}+c_{\mathrm{X}}).
  \label{eq:knowledge-acquisition-human-expert-routing-cost}
\end{equation}
式~\eqref{eq:knowledge-acquisition-human-expert-routing-cost}分别计入普通判断、专家判断和转交，不把一个转交对象误算为只有一次判断。第三种方案让专家先制定规范并培训普通人员，期待更多样本可以按规则完成。令$F_{\mathrm{D}}$、$F_{\mathrm{T}}$为开发与培训固定费，培训后仍需转交比例为$q'$，其费用为
\begin{equation}
  C_{\mathrm{train}}
  =F_{\mathrm{D}}+F_{\mathrm{T}}+nc_{\mathrm{O}}
   +q'n(c_{\mathrm{E}}+c_{\mathrm{X}}).
  \label{eq:knowledge-acquisition-human-expert-training-cost}
\end{equation}
培训增加固定费，只有转交减少等后续收益足以弥补它时才可能省钱。下面用教学数值核算，不把$q'$的减少当成实测结论。假设三方案都让$n=100$篇文档在同一规范下得到可用最终判断，普通人员不能完成的必须转交；取$c_{\mathrm{O}}=1$、$c_{\mathrm{E}}=4$、$c_{\mathrm{X}}=0.5$，全程专家$F_{\mathrm{E}}=10$，普通后转交$F_{\mathrm{R}}=20$、$q=0.20$，培训$F_{\mathrm{D}}+F_{\mathrm{T}}=70$、$q'=0.05$。代入两式与全程专家费用，依次得到
\[
  C_{\mathrm{all}}=410,\qquad
  C_{\mathrm{route}}=210,\qquad
  C_{\mathrm{train}}=192.5.
\]
再核算无法由费用代表的专家时间。每例12分钟时，全程专家用20小时，转交方案用4小时；开发培训另用6小时，培训方案则用$6+5\times12/60=7$小时。于是费用最低的培训方案并非专家时间最少的方案。所有数值均为教学假设；培训能否把转交从$20\%$降到$5\%$仍需试标，同一完成条件下的质量也须另查，不能由费用较低推出。

一般地，预计专家单例工时$t_{\mathrm{E}}$应与容量$H_{\mathrm{E}}$同单位，路由须满足$qn t_{\mathrm{E}}\leq H_{\mathrm{E}}$；开发与培训占$h_{\mathrm{D}}$、$h_{\mathrm{T}}$时，还须满足$h_{\mathrm{D}}+h_{\mathrm{T}}+q'nt_{\mathrm{E}}\leq H_{\mathrm{E}}$。转交门槛过低可能积压，过高可能迫使普通人员猜测。专家解决一例与写一条可复用规则，潜在影响范围不同；但规则覆盖多少输入仍需试标，不能先把开发费摊到全部样本就假定人人受益。

普通人员与专家也可能共享错误任务定义，因此专家身份不能替代证据，付费更高也不直接等于新增知识更多。应依据当前任务的增量价值决定何时使用专家。期限紧张时，还要比较等待专家会耽误什么、先处理其他可答文档又能完成什么，不能只优化账面费用。

\subsection{不同标注粒度与轮次的预算分配}
\label{sec:knowledge-acquisition-human-granularity-round-budget}

如果粗判断可以先排除大量无关文档，就可能只对剩余报告定位精细证据。两阶段方案先取得粗标签，再根据已知范围选择值得细化的对象；省下的是未进入第二轮的细标，但会新增重开材料和恢复上下文的工作。设第一轮$n$个对象，粗标单价$c_{\mathrm{coarse}}$；第二轮$m\leq n$个，每例重开费$c_{\mathrm{reopen}}$、细标费$c_{\mathrm{fine}}$，两轮固定费$F_1,F_2$包含各自配置、调度和所需更新，统一费用为
\begin{equation}
  C_{\mathrm{two}}
  =F_1+F_2+nc_{\mathrm{coarse}}
   +m(c_{\mathrm{reopen}}+c_{\mathrm{fine}}).
  \label{eq:knowledge-acquisition-human-two-stage-cost}
\end{equation}
式~\eqref{eq:knowledge-acquisition-human-two-stage-cost}同时记录粗标、细标和重开，不把第二轮当成只填写新增字段。一次完成全部细标签的费用为$C_{\mathrm{full}}=F_{\mathrm{full}}+nc_{\mathrm{full}}$；应在同一目标输出要求下比较。第一轮确实排除许多无须细化对象，且重复阅读、对象对齐和更新没有吞掉节省，两阶段才可能费用更低。

每轮选多少文档，也影响成本和反馈时间。大批量摊薄启动、调度和重训，却要等整批完成才利用新标签；小批量更早更新依据，但增加启动与计算。每轮应记录投入、对象覆盖、弃权及原因、独立质量证据和下游效果。预算耗尽、覆盖条件满足、验证指标达标和新增收益不足是四种不同停止理由，所需证据不同；仍未解决的重要区域要保留状态与去向。模型变得自信只说明自身分布更集中，不能单独支持停止。

表~\ref{tab:knowledge-acquisition-human-fixed-budget}改用固定120个费用单位比较配置，区别于前面要求100篇都取得最终判断的比较。设普通完整判断1单位、专家4单位、粗判断0.5单位、局部细化2单位，专家每例12分钟；前三方案准备费20，两阶段两轮配置费30。为演示核算，教学场景设$c_{\mathrm{X}}=c_{\mathrm{reopen}}=0$，转交行总费$20+80\times1+5\times(4+0)=120$，两阶段行总费$30+100\times0.5+20\times(0+2)=120$。表中对象覆盖与所得粒度不同，同额费用不表示同等质量或信息价值。

\begin{table*}[htbp]
  \centering
  \small
  \begingroup
  \setlength{\tabcolsep}{3pt}
  \renewcommand{\arraystretch}{1.15}
  \begin{tabularx}{\linewidth}{@{}>{\RaggedRight\setlength{\parindent}{0pt}\arraybackslash}p{22mm}>{\RaggedRight\setlength{\parindent}{0pt}\arraybackslash}p{14mm}>{\RaggedRight\setlength{\parindent}{0pt}\arraybackslash}p{28mm}>{\RaggedRight\setlength{\parindent}{0pt}\arraybackslash}p{16mm}>{\RaggedRight\setlength{\parindent}{0pt}\arraybackslash}p{14mm}>{\RaggedRight\setlength{\parindent}{0pt}\arraybackslash}p{32mm}>{\RaggedRight\setlength{\parindent}{0pt}\arraybackslash}X@{}}
    \toprule
    方案 & 对象数 & 判断配置 & 专家工时 & 准备费 & 逐例费用 & 待核验 \\
    \midrule
    扩大新增覆盖 & 100 & 每例1次普通判断 & 0小时 & 20 & $100\times1=100$ & 歧义与稀有类别审计 \\
    增加重复判断 & 50 & 每例2次独立普通判断 & 0小时 & 20 & $50\times2\times1=100$ & 覆盖较窄；检查共性规范错误 \\
    普通后转专家 & 80 & 每例1次普通判断；其中5例追加专家 & 1小时 & 20 & $80\times1+5\times4=100$ & 转交门槛、专家增量价值 \\
    粗标后细化 & 100 & 每例1次粗标；其中20例再细化 & 0小时 & 30 & $100\times0.5+20\times2=90$ & 80例无细标签；对象对应与重开成本 \\
    \bottomrule
  \end{tabularx}
  \endgroup
  \caption{固定资源下的人工标注方案比较。对象数为新增覆盖数，不是判断人次；准备费加逐例费均为120个费用单位，专家工时另计。四行所得覆盖、重复证据和细粒度信息不同；单价及数值为教学假设，末列保留待核验需求，不预填准确率。}
  \label{tab:knowledge-acquisition-human-fixed-budget}
\end{table*}

比较完这些配置，可以重新连接三项决策。文档是否值得获取标签，由第~\ref{sec:knowledge-acquisition-human-sample-selection}节提供依据；谁具备作答条件，由第~\ref{sec:knowledge-acquisition-human-annotator-selection}节决定；请一位专家、两位普通人员，还是用相同费用增加对象，则是本节的投入选择。三项需要联合决定，却应分别保留依据，这样才能知道失败来自选错对象、人员不匹配还是投入不足，也便于费用或容量改变后重配。

\section{本章小结}
\label{sec:knowledge-acquisition-human-summary}

人工资源有限时，我们不是单独购买更多标签，而是配置一套有明确含义的获取过程。任务与规范使不同人回答同一个问题，交互保存实际取得的答案或约束，样本与人员匹配让值得询问的材料有机会得到可用判断，预算再平衡覆盖、重复、专家与轮次。可比较方案应固定完成条件，同时满足费用和工时；交付期限还需核查依赖、并行及等待，不能由总工时直接推出。

回到故障报告，我们可以依据试标把规则清楚的常见文档交给受训人员，以分层随机样本保留覆盖，再用不确定性与代表性选补充材料。需要判读专业报告的转给专家，高风险分歧按上限追加独立判断，并另留不参与选样和调参的审计样本。缺失附件应先补材料，而不是增加作答者。粗类别若能排除大量无关对象，再定位剩余证据可能节省费用；若阅读全文才是主要成本，两轮打开可能反而更贵。

这些安排成立的条件都要由当前任务证据支持，不能直接作为通用配置。查询高分可能追逐噪声，委员会可能一致犯错，建议可能影响人工判断，重复可能复制规范误解；完成操作也不表示答案已核验。因此交付候选知识时，应能追溯问了什么、谁依据什么作答、怎样操作、哪些内容仍未知和付出了多少。后续再把多份人工答案及其他来源按共同语义融合，用独立证据评估与修正质量，而不把获取数量或模型自信当作质量结论。
